% ==============================================================================
% Figura: Esquemático Electrónico e Interconexión Esencial
% Sistema: Péndulo Invertido Rotacional de Furuta (STM32 Nucleo-F446RE)
% Construido por restricciones geométricas relativas (positioning, calc, fit)
% ==============================================================================
\begin{figure*}[!t]
\centering
% Paleta técnica local
\definecolor{mcuBorder}{HTML}{1B3A6B}
\definecolor{mcuFill}{HTML}{F2F6FC}
\definecolor{mcuHead}{HTML}{1B3A6B}
\definecolor{pwrBorder}{HTML}{9E2A2B}
\definecolor{pwrFill}{HTML}{FDF4F4}
\definecolor{pwrHead}{HTML}{9E2A2B}
\definecolor{motBorder}{HTML}{B45309}
\definecolor{motFill}{HTML}{FFFBF0}
\definecolor{motHead}{HTML}{B45309}
\definecolor{senBorder}{HTML}{0F766E}
\definecolor{senFill}{HTML}{F0FDFA}
\definecolor{senHead}{HTML}{0F766E}
\definecolor{telBorder}{HTML}{3730A3}
\definecolor{telFill}{HTML}{F1F3FF}
\definecolor{telHead}{HTML}{3730A3}
\definecolor{wirePwr}{HTML}{B91C1C}
\definecolor{wireVcc}{HTML}{DC2626}
\definecolor{wireGnd}{HTML}{1F2937}
\definecolor{wirePwm}{HTML}{C2410C}
\definecolor{wireGpio}{HTML}{1D4ED8}
\definecolor{wireScl}{HTML}{0F766E}
\definecolor{wireSda}{HTML}{15803D}
\definecolor{wireUart}{HTML}{3730A3}
\definecolor{wireMech}{HTML}{78350F}
%
\resizebox{0.94\textwidth}{!}{%
\begin{circuitikz}[
    >=Stealth,
    font=\sffamily,
    % --- Parámetros de ritmo geométrico (fuente única de verdad) ---
    /tikz/hdrnode/.style={
        rectangle,
        minimum width=#1,
        text width=\dimexpr #1 - 14pt\relax,
        align=center,
        text=white,
        font=\normalsize\bfseries,
        inner xsep=7pt,
        inner ysep=6pt
    },
    pinlbl/.style={font=\small\ttfamily\bfseries, inner sep=1pt},
    wirelbl/.style={font=\small\sffamily\bfseries, inner sep=1.5pt},
    bushead/.style={
        font=\footnotesize\sffamily\bfseries, fill=white,
        draw=#1!65, text=#1, rounded corners=2pt, inner xsep=5pt, inner ysep=2pt
    },
    pinport/.style={circle, fill=white, draw=black!80, line width=0.75pt, inner sep=1.5pt},
    solder/.style={circle, fill=black!85, inner sep=1.5pt}
]
% Dimensiones paramétricas globales
\def\mcuW{7.0cm}
\def\leftW{3.7cm}
\def\senW{3.4cm}
\def\uiW{3.5cm}
\def\pcW{3.4cm}
\def\gapLeft{3.4cm}
\def\gapRight{4.2cm}
\def\rowGap{1.00cm}
\def\pinPitch{0.78cm}
\def\botGap{1.05cm}
\def\botH{2.15cm}

% ==============================================================================
% 1. NODO RAÍZ: ENCABEZADO STM32 NUCLEO-F446RE
% ==============================================================================
\node[hdrnode=\mcuW, anchor=north] (mcu_hdr) at (0,0) {STM32 Nucleo-F446RE};

% ==============================================================================
% 2. COLUMNA DERECHA: SENSORES AS5600 BASE (I2C1) Y DISTAL (I2C2)
% ==============================================================================
\node[hdrnode=\senW, right=\gapRight of mcu_hdr.north east, anchor=north west] (as1_hdr) {AS5600 Base};

\coordinate (as1_scl) at ($(as1_hdr.south west)+(0,-0.68cm)$);
\coordinate (as1_vcc) at ($(as1_scl)+(0,-\pinPitch)$);
\coordinate (as1_sda) at ($(as1_vcc)+(0,-\pinPitch)$);
\coordinate (as1_gnd) at ($(as1_sda)+(0,-\pinPitch)$);
\coordinate (as1_bot) at ($(as1_gnd)+(0,-0.46cm)$);
\node[draw=none, fill=none, inner sep=0pt, fit=(as1_hdr)(as1_bot)] (as1) {};

\node[pinlbl, anchor=west, text=wireScl] at ($(as1_scl)+(6pt,0)$) {SCL};
\node[pinlbl, anchor=west, text=wireVcc] at ($(as1_vcc)+(6pt,0)$) {VCC};
\node[pinlbl, anchor=west, text=wireSda] at ($(as1_sda)+(6pt,0)$) {SDA};
\node[pinlbl, anchor=west, text=wireGnd] at ($(as1_gnd)+(6pt,0)$) {GND};

% Sensor 2 debajo de Sensor 1
\node[hdrnode=\senW, below=\rowGap of as1.south west, anchor=north west] (as2_hdr) {AS5600 Distal};

\coordinate (as2_scl) at ($(as2_hdr.south west)+(0,-0.68cm)$);
\coordinate (as2_vcc) at ($(as2_scl)+(0,-\pinPitch)$);
\coordinate (as2_sda) at ($(as2_vcc)+(0,-\pinPitch)$);
\coordinate (as2_gnd) at ($(as2_sda)+(0,-\pinPitch)$);
\coordinate (as2_bot) at ($(as2_gnd)+(0,-0.46cm)$);
\node[draw=none, fill=none, inner sep=0pt, fit=(as2_hdr)(as2_bot)] (as2) {};

\node[pinlbl, anchor=west, text=wireScl] at ($(as2_scl)+(6pt,0)$) {SCL};
\node[pinlbl, anchor=west, text=wireVcc] at ($(as2_vcc)+(6pt,0)$) {VCC};
\node[pinlbl, anchor=west, text=wireSda] at ($(as2_sda)+(6pt,0)$) {SDA};
\node[pinlbl, anchor=west, text=wireGnd] at ($(as2_gnd)+(6pt,0)$) {GND};

% ==============================================================================
% 3. COLUMNA IZQUIERDA: MOTOR DC, DRIVER L298N Y FUENTE 12 V
% ==============================================================================
\node[hdrnode=\leftW, left=\gapLeft of mcu_hdr.north west, anchor=north east] (motor_hdr) {Motor DC};

\node[elmech, scale=0.68, fill=white, draw=motBorder, line width=1.0pt,
      below=0.38cm of motor_hdr.south] (m_sym) {M};
\coordinate (motor_bot) at ($(m_sym.south)+(0,-0.46cm)$);
\node[draw=none, fill=none, inner sep=0pt, fit=(motor_hdr)(motor_bot)] (motor) {};

\node[font=\small\bfseries, text=motHead, anchor=south east] at ($(m_sym.left)+(-3pt,2pt)$) {$+$};
\node[font=\small\bfseries, text=motHead, anchor=south west] at ($(m_sym.right)+(3pt,2pt)$) {$-$};

\coordinate (mot_out1) at ($(motor.south west)!0.26!(motor.south east)$);
\coordinate (mot_out2) at ($(motor.south west)!0.74!(motor.south east)$);
\draw[motBorder, line width=1.05pt] (m_sym.left)  -| (mot_out1);
\draw[motBorder, line width=1.05pt] (m_sym.right) -| (mot_out2);

% Driver L298N
\node[hdrnode=\leftW, below=\rowGap of motor.south west, anchor=north west] (l298_hdr) {Driver L298N};

\coordinate (l298_ena) at ($(l298_hdr.south east)+(0,-0.56cm)$);
\coordinate (l298_in1) at ($(l298_ena)+(0,-\pinPitch)$);
\coordinate (l298_in2) at ($(l298_in1)+(0,-\pinPitch)$);
\coordinate (l298_gnd) at ($(l298_in2)+(0,-\pinPitch)$);
\coordinate (l298_bot) at ($(l298_gnd)+(0,-0.64cm)$);
\node[draw=none, fill=none, inner sep=0pt, fit=(l298_hdr)(l298_bot)] (l298n) {};

\coordinate (l298_out1) at (l298n.north -| mot_out1);
\coordinate (l298_out2) at (l298n.north -| mot_out2);
\draw[wirePwr, line width=1.35pt, ->] (l298_out1) -- (mot_out1);
\draw[wirePwr, line width=1.35pt, ->] (l298_out2) -- (mot_out2);
\node[wirelbl, text=wirePwr, anchor=east] at ($(l298_out1)!0.5!(mot_out1)+(-4pt,0)$) {OUT1};
\node[wirelbl, text=wirePwr, anchor=west] at ($(l298_out2)!0.5!(mot_out2)+( 4pt,0)$) {OUT2};

\node[pinlbl, anchor=east, text=wirePwm]  at ($(l298_ena)+(-6pt,0)$) {ENA};
\node[pinlbl, anchor=east, text=wireGpio] at ($(l298_in1)+(-6pt,0)$) {IN1};
\node[pinlbl, anchor=east, text=wireGpio] at ($(l298_in2)+(-6pt,0)$) {IN2};
\node[pinlbl, anchor=east, text=wireGnd]  at ($(l298_gnd)+(-6pt,0)$) {GND};

\coordinate (l298_vms)  at ($(l298n.south west)!0.18!(l298n.south east)$);
\coordinate (l298_pgnd) at ($(l298n.south west)!0.84!(l298n.south east)$);
\node[pinlbl, text=wirePwr, anchor=south] at ($(l298_vms)+(0,4pt)$)  {VMS};
\node[pinlbl, text=wireGnd, anchor=south] at ($(l298_pgnd)+(0,4pt)$) {GND};

% Fuente 12 V
\coordinate (psu_nw)      at ($(l298n.south west)+(0,-1.10cm)$);
\coordinate (psu_ne)      at ($(l298n.south east)+(0,-1.10cm)$);
\coordinate (psu_in_vms)  at (l298_vms |- psu_nw);
\coordinate (psu_in_gnd)  at (l298_pgnd |- psu_nw);
\coordinate (psu_src_top) at ($(psu_in_vms)+(0,-0.40cm)$);
\coordinate (psu_src_bot) at ($(psu_src_top)+(0,-1.65cm)$);
\coordinate (psu_gnd_cor) at (psu_in_gnd |- psu_src_bot);
\coordinate (psu_ne_bot)  at (psu_ne |- psu_src_bot);
\coordinate (psu_se)      at ($(psu_ne_bot)+(0,-0.40cm)$);
\node[draw=none, fill=none, inner sep=0pt, fit=(psu_nw)(psu_se)] (psu) {};

\draw[wirePwr, line width=1.35pt] (l298_vms) -- (psu_in_vms);
\node[wirelbl, text=wirePwr, anchor=east]
    at ($(l298_vms)!0.50!(psu_in_vms)+(-4pt,0)$) {+\qty{12}{\volt}};

\draw[wirePwr, line width=1.15pt] (psu_in_vms) -- (psu_src_top);
\draw[wirePwr, line width=1.15pt] (psu_src_top) to[vsource, american] (psu_src_bot);
\draw[wireGnd, line width=1.20pt]
    (psu_src_bot) -- (psu_gnd_cor) -- (psu_in_gnd) -- (l298_pgnd);

\coordinate (psu_lbl_x) at ($(psu_in_vms)!0.54!(psu_in_gnd)$);
\coordinate (psu_lbl_y) at ($(psu_src_top)!0.50!(psu_src_bot)$);
\node[font=\normalsize\bfseries, text=pwrHead, align=center]
    at (psu_lbl_x |- psu_lbl_y) {Fuente\\\qty{12}{\volt}};

% ==============================================================================
% 4. CAJA MCU CENTRAL Y PINES ORTOGONALES
% ==============================================================================
\coordinate (mcu_bot_y)   at ($(as2.south)+(0,-0.25cm)$);
\coordinate (mcu_bot_ref) at (mcu_hdr.south |- mcu_bot_y);
\node[draw=none, fill=none, inner sep=0pt, fit=(mcu_hdr)(mcu_bot_ref)] (mcu) {};

% Proyecciones de pines hacia bordes de MCU
\coordinate (mcu_scl1)  at (mcu.east |- as1_scl);
\coordinate (mcu_3v3)   at (mcu.east |- as1_vcc);
\coordinate (mcu_sda1)  at (mcu.east |- as1_sda);
\coordinate (mcu_gnd1)  at (mcu.east |- as1_gnd);

\coordinate (mcu_scl2)  at (mcu.east |- as2_scl);
\coordinate (mcu_ioref) at (mcu.east |- as2_vcc);
\coordinate (mcu_sda2)  at (mcu.east |- as2_sda);
\coordinate (mcu_gnd2)  at (mcu.east |- as2_gnd);

\coordinate (mcu_ena)   at (mcu.west |- l298_ena);
\coordinate (mcu_in1)   at (mcu.west |- l298_in1);
\coordinate (mcu_in2)   at (mcu.west |- l298_in2);
\coordinate (mcu_gnd0)  at (mcu.west |- l298_gnd);

% Pines MCU derecha
\node[pinlbl, anchor=east, text=wireScl] at ($(mcu_scl1)+(-6pt,0)$)  {PB8};
\node[pinlbl, anchor=east, text=wireVcc] at ($(mcu_3v3)+(-6pt,0)$)   {3V3};
\node[pinlbl, anchor=east, text=wireSda] at ($(mcu_sda1)+(-6pt,0)$)  {PB9};
\node[pinlbl, anchor=east, text=wireGnd] at ($(mcu_gnd1)+(-6pt,0)$)  {GND};

\node[pinlbl, anchor=east, text=wireScl] at ($(mcu_scl2)+(-6pt,0)$)  {PB10};
\node[pinlbl, anchor=east, text=wireVcc] at ($(mcu_ioref)+(-6pt,0)$) {IOREF};
\node[pinlbl, anchor=east, text=wireSda] at ($(mcu_sda2)+(-6pt,0)$)  {PB3};
\node[pinlbl, anchor=east, text=wireGnd] at ($(mcu_gnd2)+(-6pt,0)$)  {GND};

% Pines MCU izquierda
\node[pinlbl, anchor=west, text=wirePwm]  at ($(mcu_ena)+(6pt,0)$)  {PA8};
\node[pinlbl, anchor=west, text=wireGpio] at ($(mcu_in1)+(6pt,0)$)  {PB4};
\node[pinlbl, anchor=west, text=wireGpio] at ($(mcu_in2)+(6pt,0)$)  {PB5};
\node[pinlbl, anchor=west, text=wireGnd]  at ($(mcu_gnd0)+(6pt,0)$) {GND};

% ==============================================================================
% 5. BUSES DERECHOS (I2C1, I2C2) Y BUSES IZQUIERDOS (PWM/GPIO/MASA COMÚN)
% ==============================================================================
\node[bushead=senBorder, anchor=south] at ($(mcu_scl1)!0.52!(as1_scl)+(0,0.34cm)$)
    {$\text{I}^2\text{C}_1$};
\node[bushead=senBorder, anchor=south] at ($(mcu_scl2)!0.52!(as2_scl)+(0,0.34cm)$)
    {$\text{I}^2\text{C}_2$};

% Bus I2C1
\draw[wireScl, line width=1.10pt, ->]   (mcu_scl1) -- (as1_scl);
\draw[wireVcc, line width=1.15pt]       (mcu_3v3)  -- (as1_vcc);
\draw[wireSda, line width=1.10pt, <->]  (mcu_sda1) -- (as1_sda);
\draw[wireGnd, line width=1.15pt]       (mcu_gnd1) -- (as1_gnd);

\node[wirelbl, text=wireScl, above=2pt] at ($(mcu_scl1)!0.20!(as1_scl)$) {SCL1};
\node[wirelbl, text=wireVcc, above=2pt] at ($(mcu_3v3)!0.20!(as1_vcc)$)  {\qty{3.3}{\volt}};
\node[wirelbl, text=wireSda, above=2pt] at ($(mcu_sda1)!0.20!(as1_sda)$) {SDA1};
\node[wirelbl, text=wireGnd, above=2pt] at ($(mcu_gnd1)!0.20!(as1_gnd)$) {GND};

% Pull-ups I2C1
\coordinate (pu1_scl_t) at ($(mcu_scl1)!0.54!(as1_scl)$);
\coordinate (pu1_scl_b) at ($(mcu_3v3)!0.54!(as1_vcc)$);
\coordinate (pu1_sda_t) at ($(mcu_3v3)!0.78!(as1_vcc)$);
\coordinate (pu1_sda_b) at ($(mcu_sda1)!0.78!(as1_sda)$);
\draw[senBorder!85!black, line width=0.80pt]
    (pu1_scl_t) to[R, resistors/scale=0.40] (pu1_scl_b);
\draw[senBorder!85!black, line width=0.80pt]
    (pu1_sda_t) to[R, resistors/scale=0.40] (pu1_sda_b);
\node[solder, fill=wireScl] at (pu1_scl_t) {};
\node[solder, fill=wireVcc] at (pu1_scl_b) {};
\node[solder, fill=wireVcc] at (pu1_sda_t) {};
\node[solder, fill=wireSda] at (pu1_sda_b) {};
\node[font=\footnotesize, text=senHead, anchor=east] at ($(pu1_scl_t)!0.5!(pu1_scl_b)+(-5pt,0)$) {$R_{\text{pu}}$};
\node[font=\footnotesize, text=senHead, anchor=west] at ($(pu1_sda_t)!0.5!(pu1_sda_b)+( 5pt,0)$) {$R_{\text{pu}}$};

% Bus I2C2
\draw[wireScl, line width=1.10pt, ->]   (mcu_scl2)  -- (as2_scl);
\draw[wireVcc, line width=1.15pt]       (mcu_ioref) -- (as2_vcc);
\draw[wireSda, line width=1.10pt, <->]  (mcu_sda2)  -- (as2_sda);
\draw[wireGnd, line width=1.15pt]       (mcu_gnd2)  -- (as2_gnd);

\node[wirelbl, text=wireScl, above=2pt] at ($(mcu_scl2)!0.20!(as2_scl)$)  {SCL2};
\node[wirelbl, text=wireVcc, above=2pt] at ($(mcu_ioref)!0.20!(as2_vcc)$) {\qty{3.3}{\volt}};
\node[wirelbl, text=wireSda, above=2pt] at ($(mcu_sda2)!0.20!(as2_sda)$)  {SDA2};
\node[wirelbl, text=wireGnd, above=2pt] at ($(mcu_gnd2)!0.20!(as2_gnd)$)  {GND};

% Pull-ups I2C2
\coordinate (pu2_scl_t) at ($(mcu_scl2)!0.54!(as2_scl)$);
\coordinate (pu2_scl_b) at ($(mcu_ioref)!0.54!(as2_vcc)$);
\coordinate (pu2_sda_t) at ($(mcu_ioref)!0.78!(as2_vcc)$);
\coordinate (pu2_sda_b) at ($(mcu_sda2)!0.78!(as2_sda)$);
\draw[senBorder!85!black, line width=0.80pt]
    (pu2_scl_t) to[R, resistors/scale=0.40] (pu2_scl_b);
\draw[senBorder!85!black, line width=0.80pt]
    (pu2_sda_t) to[R, resistors/scale=0.40] (pu2_sda_b);
\node[solder, fill=wireScl] at (pu2_scl_t) {};
\node[solder, fill=wireVcc] at (pu2_scl_b) {};
\node[solder, fill=wireVcc] at (pu2_sda_t) {};
\node[solder, fill=wireSda] at (pu2_sda_b) {};
\node[font=\footnotesize, text=senHead, anchor=east] at ($(pu2_scl_t)!0.5!(pu2_scl_b)+(-5pt,0)$) {$R_{\text{pu}}$};
\node[font=\footnotesize, text=senHead, anchor=west] at ($(pu2_sda_t)!0.5!(pu2_sda_b)+( 5pt,0)$) {$R_{\text{pu}}$};

% Buses L298N <-> MCU
\draw[wirePwm,  line width=1.25pt, ->] (mcu_ena)  -- (l298_ena);
\draw[wireGpio, line width=1.10pt, ->] (mcu_in1)  -- (l298_in1);
\draw[wireGpio, line width=1.10pt, ->] (mcu_in2)  -- (l298_in2);
\draw[wireGnd,  line width=1.20pt]     (mcu_gnd0) -- (l298_gnd);

\node[wirelbl, text=wirePwm, above=2pt] at ($(mcu_ena)!0.50!(l298_ena)$)  {PWM};
\node[wirelbl, text=wireGnd, above=2pt] at ($(l298_gnd)!0.72!(mcu_gnd0)$) {GND};

% Masa Común
\coordinate (gnd_star_top) at ($(l298_gnd)!0.38!(mcu_gnd0)$);
\coordinate (gnd_psu_tap)  at ($(l298_pgnd)!0.55!(psu_in_gnd)$);
\coordinate (gnd_star_mid) at (gnd_star_top |- gnd_psu_tap);

\draw[wireGnd, line width=1.20pt] (gnd_star_top) -- (gnd_star_mid) -- (gnd_psu_tap);
\draw[wireGnd, line width=1.20pt] (gnd_star_mid) -- ++(0,-0.62cm) node[ground] (gnd_sym) {};
\node[solder, fill=wireGnd] at (gnd_star_top) {};
\node[solder, fill=wireGnd] at (gnd_psu_tap)  {};
\node[solder, fill=wireGnd] at (gnd_star_mid) {};
\node[font=\footnotesize\bfseries, text=wireGnd, align=center, below=2pt]
    at (gnd_sym.south) {GND común};

% ==============================================================================
% 6. CORREDOR SUPERIOR: ACOPLE MECÁNICO 80:1
% ==============================================================================
\coordinate (mech_start) at (motor.north);
\coordinate (mech_end)   at (as1.north);
\coordinate (mech_top_l) at ($(mech_start)+(0,0.80cm)$);
\coordinate (mech_top_r) at ($(mech_end)+(0,0.80cm)$);

\draw[wireMech, line width=1.35pt, densely dashed, ->, rounded corners=4pt]
    (mech_start) -- (mech_top_l) -- (mech_top_r) -- (mech_end);
\node[
    rectangle, rounded corners=2.5pt,
    draw=wireMech!75, fill=motFill, line width=0.85pt,
    font=\small\bfseries, text=wireMech, inner xsep=8pt, inner ysep=3pt
] at ($(mech_top_l)!0.50!(mech_top_r)$) {Acople 80:1};

% ==============================================================================
% 7. FILA INFERIOR: INTERFAZ NUCLEO (B1, LD2) Y PC (MATLAB)
% ==============================================================================
% Interfaz Nucleo debajo de la mitad izquierda de STM32
\coordinate (uiblk_nw) at ($(mcu.south west)+(0,-\botGap)$);
\node[draw=none, fill=none, minimum width=\uiW, minimum height=\botH,
      inner sep=0pt, anchor=north west] (uiblk) at (uiblk_nw) {};

\node[hdrnode=\uiW, anchor=south west] (uiblk_ftr) at (uiblk.south west) {Interfaz Nucleo};

\coordinate (uiblk_b1_top)  at ($(uiblk.north west)!0.28!(uiblk.north east)$);
\coordinate (uiblk_ld2_top) at ($(uiblk.north west)!0.58!(uiblk.north east)$);

\coordinate (mcu_pc13) at (mcu.south -| uiblk_b1_top);
\coordinate (mcu_pa5)  at (mcu.south -| uiblk_ld2_top);

\node[pinlbl, text=telBorder, anchor=south] at ($(mcu_pc13)+(0,5pt)$) {PC13};
\node[pinlbl, text=telBorder, anchor=south] at ($(mcu_pa5)+(0,5pt)$)  {PA5};

\coordinate (b1_s1) at ($(uiblk_b1_top)+(0,-0.24cm)$);
\coordinate (b1_s2) at ($(uiblk_b1_top)+(0,-1.06cm)$);
\draw[telBorder, line width=1.10pt, <-] (mcu_pc13) -- (uiblk_b1_top);
\draw[telBorder, line width=0.95pt]     (uiblk_b1_top) -- (b1_s1)
    to[push button, switches/scale=0.52] (b1_s2) node[ground, scale=0.55] {};
\node[font=\small\bfseries, text=telHead, anchor=east]
    at ($(b1_s1)!0.48!(b1_s2)+(-9pt,0)$) {B1};

\coordinate (ld2_s1) at ($(uiblk_ld2_top)+(0,-0.24cm)$);
\coordinate (ld2_s2) at ($(uiblk_ld2_top)+(0,-1.06cm)$);
\draw[telBorder, line width=1.10pt, ->] (mcu_pa5) -- (uiblk_ld2_top);
\draw[telBorder, line width=0.95pt]     (uiblk_ld2_top) -- (ld2_s1)
    to[leD, diodes/scale=0.50, color=wireSda] (ld2_s2) node[ground, scale=0.55] {};
\node[font=\small\bfseries, text=wireSda, anchor=west]
    at ($(ld2_s1)!0.48!(ld2_s2)+(17pt,0)$) {LD2};

% PC (MATLAB) debajo de AS5600 Distal (alineado en fila inferior)
\coordinate (pc_ne) at (as2.east |- uiblk.north);
\node[draw=none, fill=none, minimum width=\pcW, minimum height=\botH,
      inner sep=0pt, anchor=north east] (pc) at (pc_ne) {};
\node[font=\normalsize\bfseries, text=telHead, align=center] at (pc.center) {PC\\(MATLAB)};

% Conexiones UART PA2 (TX) y PA3 (RX)
\coordinate (mcu_pa3) at ($(mcu.south west)!0.59!(mcu.south east)$);
\coordinate (mcu_pa2) at ($(mcu.south west)!0.87!(mcu.south east)$);
\coordinate (pc_tx_in)  at ($(pc.north west)!0.36!(pc.south west)$);
\coordinate (pc_rx_out) at ($(pc.north west)!0.74!(pc.south west)$);
\coordinate (uart_turn_tx) at (mcu_pa2 |- pc_tx_in);
\coordinate (uart_turn_rx) at (mcu_pa3 |- pc_rx_out);

\node[pinlbl, text=wireUart, anchor=south] at ($(mcu_pa3)+(0,5pt)$) {PA3 (RX)};
\node[pinlbl, text=wireUart, anchor=south] at ($(mcu_pa2)+(0,5pt)$) {PA2 (TX)};

\draw[wireUart, line width=1.15pt, ->, rounded corners=3pt]
    (mcu_pa2) -- (uart_turn_tx) -- (pc_tx_in);
\draw[wireUart, line width=1.15pt, <-, rounded corners=3pt]
    (mcu_pa3) -- (uart_turn_rx) -- (pc_rx_out);

\node[wirelbl, text=wireUart, above=2pt] at ($(uart_turn_tx)!0.50!(pc_tx_in)$) {TX};
\node[wirelbl, text=wireUart, above=2pt] at ($(uart_turn_rx)!0.55!(pc_rx_out)$) {RX};
\node[bushead=telBorder, anchor=south] at ($(uart_turn_tx)!0.50!(pc_tx_in)+(0,0.56cm)$) {UART};

% ==============================================================================
% 8. CAPA DE FONDO Y MARCOS
% ==============================================================================
\begin{scope}[on background layer]
    % MCU
    \path[fill=mcuFill, rounded corners=4pt] (mcu.north west) rectangle (mcu.south east);
    \coordinate (mcu_hb_w) at (mcu.west |- mcu_hdr.south);
    \coordinate (mcu_hb_e) at (mcu.east |- mcu_hdr.south);
    \fill[mcuHead, rounded corners=4pt] (mcu.north west) rectangle (mcu_hb_e);
    \fill[mcuHead] ($(mcu_hb_w)+(0,5pt)$) rectangle (mcu_hb_e);
    \draw[mcuBorder, rounded corners=4pt, line width=1.4pt] (mcu.north west) rectangle (mcu.south east);

    % Sensores AS5600 Base y Distal
    \foreach \s in {as1, as2} {
        \path[fill=senFill, rounded corners=3.5pt] (\s.north west) rectangle (\s.south east);
        \coordinate (\s_hb_w) at (\s.west |- \s_hdr.south);
        \coordinate (\s_hb_e) at (\s.east |- \s_hdr.south);
        \fill[senHead, rounded corners=3.5pt] (\s.north west) rectangle (\s_hb_e);
        \fill[senHead] ($(\s_hb_w)+(0,5pt)$) rectangle (\s_hb_e);
        \draw[senBorder, rounded corners=3.5pt, line width=1.2pt] (\s.north west) rectangle (\s.south east);
    }

    % Motor DC
    \path[fill=motFill, rounded corners=3.5pt] (motor.north west) rectangle (motor.south east);
    \coordinate (mot_hb_w) at (motor.west |- motor_hdr.south);
    \coordinate (mot_hb_e) at (motor.east |- motor_hdr.south);
    \fill[motHead, rounded corners=3.5pt] (motor.north west) rectangle (mot_hb_e);
    \fill[motHead] ($(mot_hb_w)+(0,5pt)$) rectangle (mot_hb_e);
    \draw[motBorder, rounded corners=3.5pt, line width=1.2pt] (motor.north west) rectangle (motor.south east);

    % Driver L298N
    \path[fill=pwrFill, rounded corners=3.5pt] (l298n.north west) rectangle (l298n.south east);
    \coordinate (l298_hb_w) at (l298n.west |- l298_hdr.south);
    \coordinate (l298_hb_e) at (l298n.east |- l298_hdr.south);
    \fill[pwrHead, rounded corners=3.5pt] (l298n.north west) rectangle (l298_hb_e);
    \fill[pwrHead] ($(l298_hb_w)+(0,5pt)$) rectangle (l298_hb_e);
    \draw[pwrBorder, rounded corners=3.5pt, line width=1.2pt] (l298n.north west) rectangle (l298n.south east);

    % Fuente 12 V
    \path[fill=pwrFill!65!white, rounded corners=3.5pt] (psu.north west) rectangle (psu.south east);
    \draw[pwrBorder, dashed, rounded corners=3.5pt, line width=1.05pt] (psu.north west) rectangle (psu.south east);

    % Interfaz Nucleo
    \path[fill=telFill!60!white, rounded corners=3.5pt] (uiblk.north west) rectangle (uiblk.south east);
    \coordinate (ui_ft_w) at (uiblk.west |- uiblk_ftr.north);
    \coordinate (ui_ft_e) at (uiblk.east |- uiblk_ftr.north);
    \fill[telHead!88!black, rounded corners=3.5pt] (ui_ft_w) rectangle (uiblk.south east);
    \fill[telHead!88!black] (ui_ft_w) rectangle ($(ui_ft_e)+(0,-5pt)$);
    \draw[telBorder!85!black, rounded corners=3.5pt, line width=1.05pt] (uiblk.north west) rectangle (uiblk.south east);

    % PC (MATLAB)
    \path[fill=telFill, rounded corners=3.5pt] (pc.north west) rectangle (pc.south east);
    \draw[telBorder, rounded corners=3.5pt, line width=1.2pt] (pc.north west) rectangle (pc.south east);
\end{scope}

% Puertos de conexión circulares sobre los bordes
\foreach \p in {
    mcu_scl1, mcu_3v3, mcu_sda1, mcu_gnd1, as1_scl, as1_vcc, as1_sda, as1_gnd,
    mcu_scl2, mcu_ioref, mcu_sda2, mcu_gnd2, as2_scl, as2_vcc, as2_sda, as2_gnd,
    mot_out1, mot_out2, l298_out1, l298_out2, l298_vms, l298_pgnd, psu_in_vms, psu_in_gnd,
    l298_ena, l298_in1, l298_in2, l298_gnd, mcu_ena, mcu_in1, mcu_in2, mcu_gnd0,
    mcu_pc13, mcu_pa5, mcu_pa3, mcu_pa2, uiblk_b1_top, uiblk_ld2_top, pc_tx_in, pc_rx_out} {
    \node[pinport] at (\p) {};
}

\end{circuitikz}%
}%
\caption{Esquemático electrónico simplificado del sistema mecatrónico sobre la placa STM32 Nucleo-F446RE: etapa de potencia \qty{12}{\volt} y puente H L298N (rojo: potencia; naranja: PWM \texttt{PA8}; azul: dirección \texttt{PB4}/\texttt{PB5}; negro: masa común), doble bus Soft-$\text{I}^2\text{C}$ con resistencias \textit{pull-up} $R_{\text{pu}}$ hacia los encoders AS5600 base y distal (verde azulado/verde: \texttt{SCL}/\texttt{SDA}; carmesí: \qty{3.3}{\volt}; línea punteada: acople mecánico $80:1$), interfaz local (\texttt{B1}, \texttt{LD2}) y enlace \texttt{UART} (\texttt{PA2}/\texttt{PA3}) hacia PC (MATLAB).}
\label{fig:esquematico_electronico}
\end{figure*}
